\documentclass{article}

\usepackage[T1]{fontenc}
\usepackage[french]{babel}
\usepackage{graphicx}
\usepackage{hyperref}
\usepackage{listings}
\usepackage{fancyhdr}

\newcommand{\latex}{\LaTeX}

% -----------page 1-----------
\title{Programmation FPGA avec Eclat : performances}
\author {Dahmane Ilan \\ 
APR/LIP6 et IRILL   
}
\date{1 Juin 2026 - 31 Juillet 2026}

\begin{document}
\maketitle

\vfill

\begin{center}
\begin{minipage}{0.4\textwidth}
\textbf{Tuteur universitaire :}\\
Eugene Asarin\\
Giulio Manzonetto
\end{minipage}
\hfill
\begin{minipage}{0.4\textwidth}
\textbf{Tuteur de stage :}\\
Loïc Sylvestre\\
Emmanuel Chailloux
\end{minipage}
\end{center}

% -----------page 2-----------
\newpage
\renewcommand{\contentsname}{Sommaire}
\tableofcontents


% -----------page 3-----------
\newpage
\section{Introduction}
\paragraph{} Dans l'industrie informatique, les processeurs généralistes (CPU) constituent 
depuis plusieurs décennies la solution dominante pour l'exécution des programmes. 
Leur architecture est conçue pour offrir une grande flexibilité et pouvoir traiter une grande variété de tâches. 
Cependant, cette généralité limite leur efficacité lorsqu'il s'agit d'exploiter 
le parallélisme massif présent dans certaines applications.

Pour répondre à certains besoins spécifiques, les GPU (Graphics Processing Units) ont été développés 
initialement pour accélérer le rendu graphique. 
Leur architecture massivement parallèle les rend particulièrement efficaces pour 
les applications reposant sur de nombreuses opérations similaires appliquées sur 
de grands ensembles de données, comme le traitement d'image ou l'intelligence artificielle.

Cependant, toutes les applications ne se prêtent pas naturellement à ce modèle de calcul. 
Certaines nécessitent des architectures spécialisées capables d'exploiter un parallélisme différent, 
avec des contraintes fortes en termes de latence, de consommation énergétique ou de traitement temps réel. 
C'est dans ce contexte que les FPGA (Field Programmable Gate Array) constituent 
une solution intéressante grâce à leur capacité de reconfiguration \cite{sentieys:hal-00716772} et 
leur aptitude à implémenter des architectures matérielles spécialisées.

Cependant, leur programmation reste un domaine complexe nécessitant traditionnellement 
une description matérielle bas niveau. Le langage Eclat \cite{these_eclat} s'inscrit dans cette problématique 
en proposant une approche permettant de décrire des calculs matériels avec un niveau d'abstraction supérieur 
tout en conservant une traduction efficace vers des architectures FPGA. 

Si les langages de plus haut niveau comme Eclat permettent de simplifier la programmation des FPGA, 
ils soulèvent néanmoins une question importante. 
Les concepteurs matériels raisonnent traditionnellement en termes de registres, de portes logiques, 
de mémoires et de cycles d'horloge, tandis que les programmeurs sont habitués à 
manipuler des abstractions de plus haut niveau, des structures de données complexes et 
différents paradigmes de programmation. Il devient alors nécessaire de déterminer 
quels styles de programmation permettent de concilier expressivité, sûreté et efficacité du matériel généré.

Dans quelle mesure les styles de programmation fonctionnel et impératif proposés par Eclat 
permettent-ils de décrire efficacement des algorithmes réalistes tout en générant des 
architectures matérielles performantes ?

Pour répondre à cette problématique, ce stage consiste à comparer les approches fonctionnelle et 
impérative au travers de l'implémentations d'un processeur LC-3\footnote{Le LC-3 est un processeur pédagogique conçu par 
Yale N.Patt et Sanjay J.Patel. Introduit dans leur ouvrage 
Introduction to Computing Systems afin d'enseigner les principes fondamentaux de l'architecture des ordinateurs.}, 
étudié à l'université de Paris-Cité durant le cours 
de Circuits et Architecture enseigné par le professeur Adrien Guatto. Le processeur LC-3 constitue un excellent candidat 
pour cette évaluation grâce à son compromis entre simplicité et représentativité. 
Son architecture volontairement réduite permet une compréhension complète de son fonctionnement, 
tout en mettant en œuvre de nombreux mécanismes caractéristiques des processeurs modernes. 
Il sollicite ainsi plusieurs fonctionnalités importantes d'Eclat, telles que la gestion de l'état, 
des registres, de la mémoire et des structures de données.
Cette complexité maîtrisée en fait un support pertinent pour tester un large panel des fonctionnalités 
du langage et comparer différentes approches d'implémentation. De plus, l'utilisation d'un même processeur 
comme base d'expérimentation permet de conserver un comportement fonctionnel identique entre les différentes versions, 
tout en observant l'impact du style de programmation choisi sur l'architecture matérielle générée, 
les performances obtenues et les ressources utilisées.
Les deux styles sont évalués selon différents critères tels que la 
facilité d'écriture, les performances obtenues après synthèse et l'occupation des ressources matérielles sur FPGA.

% -----------page 4-----------
\newpage
\section{FPGA et programmation}
\paragraph{} Les FPGA représentent une évolution majeure dans le domaine de la conception matérielle. 
Contrairement aux circuits intégrés spécifiques (ASIC), dont la fabrication nécessite des coûts de conception initiaux très importants, 
les FPGA permettent de configurer après fabrication une architecture matérielle adaptée à un besoin particulier. 
Cette flexibilité réduit considérablement les coûts et les délais de développement, 
notamment pour les systèmes nécessitant des évolutions fréquentes ou des productions en volumes limités.

La particularité principale des FPGA réside dans leur capacité de reconfiguration \cite{these_hls}. 
Ils permettent d'exploiter un parallélisme spatial en instanciant directement plusieurs unités de calcul matérielles.
Cette caractéristique en fait des accélérateurs matériels particulièrement adaptés à certaines applications industrielles nécessitant de fortes performances, 
comme le traitement du signal, l'intelligence artificielle, le calcul scientifique ou encore les systèmes embarqués temps réel.

Cependant, l'exploitation de ces architectures nécessite traditionnellement l'utilisation de langages de description matérielle (HDL), tels que le VHDL ou le Verilog. 
Contrairement aux langages de programmation classiques, ces langages ne décrivent pas une succession d'instructions exécutées par un processeur, 
mais une architecture matérielle composée de composants logiques fonctionnant potentiellement en parallèle. 
Cette approche offre un contrôle très fin sur l'implémentation matérielle, mais introduit également une complexité importante et 
nécessite des connaissances spécifiques en conception numérique. 
L'objectif est donc de proposer des outils permettant de rapprocher la programmation FPGA des méthodes utilisées en programmation logicielle classique, 
tout en conservant un contrôle suffisant sur l'architecture générée.. 

Un des exemples les plus connus est HLS (High-level Synthesis) afin de se concentrer sur l'aspect algorithmique et n'ont pas sur l'architecture du circuit. 
Nous pouvons notamment écrire dans un sous ensemble de C/C++ (pointeur et allocation dynamique parfois exclus) afin que l'outil de synthèse de haut niveau 
parse votre programme et le synthétise en HDL. 
Mais cette approche à une limite. Le contrôle de l'architecture par le concepteur du programme est grandement limité et est délégué à l'outil de synthèse (le parallélisme
est inférer etc... creuser un peu plus). 

C'est dans cette perspective que s'inscrit Eclat. Ce langage propose un paradigme de programmation plus proche des langages logiciels traditionnels (OCaml-like) 
tout en conservant la possibilité de générer une architecture matérielle efficace. Le compilateur Eclat permet ainsi de traduire 
les programmes écrits dans ce langage en une description matérielle VHDL synthétisable, pouvant ensuite être implémentée sur un FPGA.


% -----------page 5-----------
\newpage
\section{Éxperimentations}
\subsection{Processeur LC-3}
Pour répondre à notre question, nous allons reprendre un ancien projet de mon premier semestre de master. Ce projet portait sur la réalisation d'un 
processeur à but éducatif nommé LC-3. C'est un processeur très rudimentaire, comprenant d'environ 15 instructions de base. Ces dernières se décomposent 
en trois familles : : les instructions arithmétiques et logiques, les instructions de chargement et rangement et les instructions de 
branchement (appelées aussi instructions de saut ou encore instructions de contrôle).

Les instructions arithmétiques et logiques du LC-3 sont au nombre de trois. Le LC-3 possède une
instruction ADD pour l’addition, une instruction AND pour le et logique bit à bit et une instruction NOT
pour la négation bit à bit. Le résultat de ces trois opérations est toujours placé dans un registre. Les
deux opérandes peuvent être soit les contenus de deux registres soit le contenu d’un registre et une
constante pour ADD et AND.

Les instructions de chargement et rangement comprennent des instructions (avec des noms
commencent par LD) permettant de charger un registre avec une valeur et des instructions (avec des
noms commencent par ST) permettant de ranger en mémoire le contenu d’un registre. Ces instructions
se différencient par leurs modes d’adressage qui peut être immédiat, direct, indirect ou relatif. Les
instructions de chargement sont LD, LDI, LDR, LEA et les instructions de rangement sont ST, STI et
STR

Les instructions de branchement comprennent les deux instructions de saut BR et JMP, les deux
instructions d’appel de sous-routine JSR et JSRR, une instruction TRAP d’appel système (pas implementé), une
instruction RET de retour de sous-routine et une instruction RTI de retour d’interruption (pas implementé).

De plus, le micro-processeur LC-3 dispose de 8 registres généraux 16 bits appelés R0...R7 et comme tous les micro-processeurs, le LC-3 
dispose d’un compteur de programme PC et d’un registre d’instruction IR qui sont tous les deux des registres 16 bits.

Les indicateurs N, Z et P sont des registres 1 bit. Ils sont positionnés dès qu’une nouvelle valeur est mise dans un des registres généraux
R0...R7. Ceci a lieu lors de l’exécution d’une instruction logique (NOT et AND), arithmétique (ADD)
ou d’une instruction de chargement (LD, LDI, LDR et LEA). Ils indiquent respectivement si cette
nouvelle valeur est négative (n), nulle (z) et positive (p). Ces indicateurs sont utilisés par l’instruction
de branchement conditionnel BR pour savoir si le branchement doit être pris ou non.

La mémoire du LC-3 est organisée par mots de 16 bits. L’adressage du LC-3 est également de 16 bits.
La mémoire du LC-3 est donc formée de 216 mots de 16 bits, c’est-à-dire 128 KiB avec des adresses
de 0000 à FFFF en hexadécimal. (insérer image ?)

\subsection{Première implantation de LC3 en Eclat}
Pour la première implémentation en Eclat du processeur, j'ai choisi de privilégier une approche impérative, 
plus proche des langages de programmation classiques et plus pédagogique pour prendre en main Eclat. 
Les registres ainsi que la mémoire ont été représentés par des tableaux, qui constituent la principale structure de données mutable du langage. 
Les modifications sont effectuées en place, c'est-à-dire par réécriture directe des cases mémoire concernées, 
sans recopier l'ensemble de la structure. Le tableau représentant la mémoire contient des valeurs du type somme value, 
capable de représenter les différents éléments susceptibles d'être stockés en mémoire, tels qu'une instruction, une adresse, 
un caractère ou une constante. Cette première implémentation est visible en annexe. 

\subsection{Résultats des benchmarks}
J'ai commencé par allumer simultannémament un paquet de leds 
afin de comprendre la chaîne de compilation d'un langage de description matériel jusqu'à la production du .sof qui décrit le circuit et permet de flash sur la carte
le circuit synthétiser par eclat.

- Parler du nombre de bloc logique utilisé pour les deux synthèse
- Parler du temps en cycle 
- Parler de la facilité d'écriture 
- Insérer des benchs avec un programme équivalent mais synthétisé avec la vm OCaml en Eclat


% -----------page 6-----------
\newpage
\section{Conclusion}
\paragraph{} Nous pouvons voir que les deux approches ont leur qualités et leur défauts, d’un côté nous avons de la performance temporelle pour le style fonctionnel 
qui permet de réduire les couts d’accès du à l’abscence d’accès concurrent (qui coute 1 cycle par accès) 
mais d’un autre côté nous avons une explosion du nombre de blocs logiques programmables utilisé en raison 
des nombreuses copies de bus de fils etc… De l’autre nous avons l’exacte opposé, 
en terme de complexité pire mais qui en terme de taille de circuit et de blocs logiques programmables est bien supérieur et 
permet de faire des programme bien plus gros. Une seconde explication passe par 
la logique d’implémentation des structure de données des deux styles. 
Les tableaux sont dans la BRAM (je crois car il y a aussi la SDRAM mais c'est une mémoire externe et c'est pour les grosses donneés non ?) 
de la carte ce qui ne limite pas sa taille 
tandis que les vecteurs sont des suites de blocs logiques programmables ce qui limite sa taille à 128kb (voir photo ou plutot donner une équivalent en terme de 
nombre d'entiers 32 bits, plus clair)
hypothétique en attendant de finir les benchmarks et donc la conclusion.....

\newpage

\bibliographystyle{unsrt}     
\bibliography{bibliographie}
\end{document}